\section{Two models at the prescribed endpoint}\label{sec:endpoint}

We construct an object whose index sort is \(V_\eta\) and whose order sort
has cardinality \(\delta\), where
\[
  \delta=\omega_1,\qquad \eta=\delta^+=\omega_2.
\]
The construction assigns to every finite tuple a countable extensional
membership diagram.  Actual ranks below \(\eta\) are replaced, separately
at each order point, by tags below \(\delta\).  For any one tuple, all
required inequalities hold outside a countable set of order points.
The exception set is allowed to depend on the tuple, as in the definition
of \(K\).

\subsection{Ranks represented modulo countable sets}

\begin{lemma}\label{end:domination}
There are functions \(f_\alpha:\delta\longrightarrow\delta\), for
\(\alpha<\eta\), such that for all \(\gamma<\alpha<\eta\) the set
\[
  E_{\gamma,\alpha}
  =\{\beta<\delta:f_\gamma(\beta)\geq f_\alpha(\beta)\}
\]
is countable.
\end{lemma}

\begin{proof}
For each \(\alpha<\eta=\delta^+\), fix an injection
\(j_\alpha:\alpha\longrightarrow\delta\).  These form a set-sized
family of choices.  Define the functions by transfinite recursion,
putting
\begin{equation}\label{end:diagonal-formula}
  f_\alpha(\beta)
  =
  \sup\{f_\gamma(\beta)+1:
        \gamma<\alpha,\ j_\alpha(\gamma)<\beta\},
  \qquad \sup\varnothing=0.
\end{equation}
For each \(\beta<\omega_1\), the set of indices in this supremum
is countable, because \(j_\alpha\) is injective and \(\beta\) is a
countable ordinal.  Regularity of \(\omega_1\) implies that the
supremum is below \(\delta\).  Thus the recursion defines a
function into \(\delta\) at every stage \(\alpha<\eta\).

If \(\gamma<\alpha\) and \(\beta>j_\alpha(\gamma)\), the term
\(f_\gamma(\beta)+1\) occurs in
\eqref{end:diagonal-formula}.  Hence
\[
  f_\gamma(\beta)<f_\alpha(\beta),
  \qquad
  E_{\gamma,\alpha}\subseteq j_\alpha(\gamma)+1.
\]
The latter ordinal is countable, proving the assertion.
\end{proof}

\begin{lemma}\label{end:hierarchy-size}
The cumulative hierarchy satisfies
\[
  |V_{\omega+\gamma}|=\beth_\gamma
  \quad\text{for every ordinal }\gamma.
\]
In particular, \(|V_\eta|=\beth_\eta\), and every element of \(V_\eta\)
has rank below \(\eta\).
\end{lemma}

\begin{proof}
The set \(V_\omega\) of hereditarily finite sets is countably
infinite, giving the assertion at \(\gamma=0\).  At a successor
stage,
\[
  |V_{\omega+(\gamma+1)}|
  =|\mathcal P(V_{\omega+\gamma})|
  =2^{\beth_\gamma}
  =\beth_{\gamma+1}.
\]
If \(\gamma\) is a nonzero limit ordinal, continuity of ordinal
addition and of the cumulative hierarchy gives
\[
  V_{\omega+\gamma}
  =\bigcup_{\zeta<\gamma} V_{\omega+\zeta}.
\]
The inductive hypothesis gives a lower bound of
\(\sup_{\zeta<\gamma}\beth_\zeta=\beth_\gamma\) for the
cardinality of this union.  It also gives the upper bound
\(|\gamma|\cdot\beth_\gamma\).  Strict increase of the beth
sequence implies \(|\gamma|\leq\beth_\gamma\): the map
\(\zeta\mapsto\beth_\zeta\), for \(\zeta<\gamma\), is an injection
into the ordinal \(\beth_\gamma\).  The upper bound is therefore
\(\beth_\gamma\), completing the induction.

Since \(\omega+\omega_2=\omega_2\), the first conclusion follows.
The rank assertion is the defining property of membership in
\(V_\eta\).
\end{proof}

\subsection{Compatible countable membership diagrams}

Put \(D=V_\eta\).  We first choose a countable set of diagram
entries for each nonempty finite subset of \(D\).  The closure
below ensures extensionality separately for every such set.
Starting from a nonempty countable set of entries, we repeatedly
adjoin one distinguishing membership witness for each pair of
distinct entries. The resulting set remains countable and has an
extensional membership relation, even when the entries themselves
are large sets.

\begin{lemma}\label{end:countable-closures}
There is a family
\[
  (D_X:\varnothing\ne X\in[D]^{<\omega})
\]
of nonempty countable subsets of \(D\) such that:
\begin{enumerate}[label=\textup{(\roman*)}]
  \item \(X\subseteq D_X\);
  \item \(D_Y\subseteq D_X\) whenever
        \(\varnothing\ne Y\subseteq X\);
  \item the relation of actual membership on \(D_X\) is
        extensional: for distinct \(a,b\in D_X\), some
        \(w\in D_X\) belongs to exactly one of \(a,b\).
\end{enumerate}
We may also fix surjections \(d_X:\omega\longrightarrow D_X\)
such that \(d_{\{x\}}(0)=x\) for each \(x\in D\).
\end{lemma}

\begin{proof}
Fix a well-order of the set \(D\).  For distinct \(a,b\in D\),
let \(w(a,b)\) be the first element of \(a\SymDiff b\) in this
well-order.  This witness exists by set extensionality and
belongs to \(D\), since \(D=V_\eta\) is transitive.

If \(A\subseteq D\) is nonempty and countable, define
\[
  A_0=A,\qquad
  A_{n+1}=A_n\cup
      \{w(a,b):a,b\in A_n,\ a\ne b\},
  \qquad
  C(A)=\bigcup_{n<\omega}A_n.
\]
Every \(A_n\), and hence \(C(A)\), is countable.  Any two
distinct elements of \(C(A)\) lie together in some \(A_n\);
their chosen witness lies in \(A_{n+1}\).  Thus membership
on \(C(A)\) is extensional.

Recursively on the finite cardinality of \(X\), set
\[
  D_X=C\left(
    X\cup
    \bigcup_{\varnothing\ne Y\subsetneq X}D_Y
  \right).
\]
The seed is nonempty and countable, because there are only
finitely many proper subsets of \(X\).  This proves
\textup{(i)} and \textup{(iii)} and directly gives
\textup{(ii)} for proper inclusions; equality gives the
remaining case.

Each \(D_X\) is nonempty and at most countable, so it admits
a surjection from \(\omega\), allowing repetitions when it
is finite.  For \(X=\{x\}\), choose one whose initial entry
is \(x\).  All these choices are indexed by a set, and
the axiom of choice supplies the family of surjections.
\end{proof}

Fix the functions from Lemma~\ref{end:domination} and the
sets and enumerations from Lemma~\ref{end:countable-closures}.
For every nonempty finite \(X\subseteq D\), define
\begin{equation}\label{end:bad-set}
  F_X=
  \bigcup_{\substack{x,y\in D_X\\y\in x}}
       E_{\rank(y),\rank(x)}.
\end{equation}
Each index pair is legitimate:
\(\rank(y)<\rank(x)<\eta\).  There are countably many
pairs in the union, so \(F_X\) is countable.  Moreover,
\begin{equation}\label{end:bad-monotone}
  \varnothing\ne Y\subseteq X
  \quad\Longrightarrow\quad
  F_Y\subseteq F_X.
\end{equation}
Indeed, \(D_Y\subseteq D_X\), so every edge used to define
\(F_Y\) is also used to define \(F_X\).

We use the label sets and the coding maps of
Section~\ref{sec:colors}.  For an ordered tuple
\[
  t=(x_1,\ldots,x_k)
\]
of distinct elements of \(D\), write \(X=\{x_1,\ldots,x_k\}\).
If \(\varnothing\ne J\subseteq[k]\), let
\(X_J=\{x_j:j\in J\}\).  Define a map from the label set
\(\mathcal L_k\) by
\[
  q_t(J,i)=d_{X_J}(i).
\]
For each \(\beta<\delta\), define an arity-\(k\) diagram
\(\mathcal D_{t,\beta}=(\mathcal E_t,R_t,\rho_{t,\beta})\) by
\begin{align}
  \ell\mathcal E_t\ell'
  &\quad\Longleftrightarrow\quad
      q_t(\ell)=q_t(\ell'), \label{end:diagram-equality}\\
  \ell R_t\ell'
  &\quad\Longleftrightarrow\quad
      q_t(\ell)\in q_t(\ell'), \label{end:diagram-membership}\\
  \rho_{t,\beta}(\ell)
  &=f_{\rank(q_t(\ell))}(\beta).
      \label{end:diagram-tags}
\end{align}
All entries of \(q_t\) lie in \(D_X\), by the inclusion
property of the closures.  Since the labels \(([k],i)\)
enumerate all of \(D_X\), the map \(q_t\) is onto \(D_X\).

\begin{lemma}\label{end:diagrams}
For every tuple \(t\) as above, the equality quotient of
\(\mathcal D_{t,\beta}\) identifies with \(D_X\), with its
actual membership relation, for every \(\beta<\delta\).
The diagrams satisfy the following additional properties.
\begin{enumerate}[label=\textup{(\roman*)}]
  \item If \(\varnothing\ne J\subseteq[k]\), restriction of
        \(\mathcal D_{t,\beta}\) to the labels belonging to
        \(J\), with the increasing reindexing of positions,
        is exactly \(\mathcal D_{t\restriction J,\beta}\).
  \item If \(\beta\notin F_X\), then every diagram
        \(\mathcal D_{t\restriction J,\beta}\), including the full
        diagram, is valid.
\end{enumerate}
Consequently, for every \(\beta\notin F_X\), the colors
\[
  s_{t\restriction J,\beta}
    =\iota_{|J|}(\mathcal D_{t\restriction J,\beta}),
  \qquad \varnothing\ne J\subseteq[k],
\]
pass the full subsequence test of
Definition~\ref{col:test}.
\end{lemma}

\begin{proof}
Surjectivity of \(q_t\), together with
\eqref{end:diagram-equality}, identifies the quotient
bijectively with \(D_X\).
Equation~\eqref{end:diagram-membership} identifies its
relation with membership on \(D_X\), which is
extensional.  The tags in \eqref{end:diagram-tags}
depend only on the image under \(q_t\), so both the
relation and the tags respect the equality relation.

For \textup{(i)}, list \(J\) increasingly as
\(\{j_1<\cdots<j_m\}\).  A label \((A,i)\) of the
subtuple \(t\restriction J=(x_{j_1},\ldots,x_{j_m})\) corresponds
to \((\{j_a:a\in A\},i)\) in the original tuple.
Both labels denote the same entry
\[
  d_{\{x_{j_a}:a\in A\}}(i).
\]
Their equality relations, membership relations, and
tags therefore agree exactly.  This proof applies to
every ordered tuple; the enumeration \(d_Y\) was fixed
for the underlying set \(Y\).

For the full diagram in \textup{(ii)}, suppose
\(\beta\notin F_X\).  Every membership edge \(y\in x\)
with \(x,y\in D_X\) then satisfies
\[
  f_{\rank(y)}(\beta)<f_{\rank(x)}(\beta)
\]
by \eqref{end:bad-set}.  This is precisely the strict
tag inequality required for that edge.  Finally,
\(q_t(\{j\},0)=x_j\), so the distinguished singleton
labels are pairwise inequivalent.  All validity
conditions now hold.

For a subtuple with underlying set \(Y=X_J\), its
quotient is \(D_Y\), whose membership relation is
extensional by its own construction.  In addition,
\(\beta\notin F_Y\) follows from
\eqref{end:bad-monotone}.  The same argument proves
the validity of this subtuple diagram.  Combined
with \textup{(i)}, these are exactly the conditions
of the full subsequence test.
\end{proof}

For all \(t\) and all \(\beta<\delta\), define
\begin{equation}\label{end:prescribed-colors}
  s_{t,\beta}=\iota_k(\mathcal D_{t,\beta}),
  \qquad k=|t|.
\end{equation}
The coding maps are defined on arbitrary diagrams, so
this prescription also makes sense at the exceptional
coordinates in \(F_X\).

\subsection{Realizing the colors in an object}

Let \(B=\delta\times\Q\) with its lexicographic order.
This is a dense linear order without endpoints:
within any rational block there are points above and
below each given point; if two points belong to
different blocks, a larger point in the earlier block
lies strictly between them.  The interval between
points in blocks \(\alpha\leq\alpha'<\omega_1\) is
contained in
\[
  \bigl([\alpha,\alpha']\cap\omega_1\bigr)\times\Q,
\]
which is countable.  Thus \(B\) meets the order
requirements for a basic structure and has
cardinality \(\delta\).  Fix a bijection
\(\beta\mapsto b_\beta\) from \(\delta\) to \(B\);
this indexing is independent of the order.

Take \(I=D\). For every nonempty finite ordered tuple \(t\) of
distinct elements of \(I\), and every \(\beta<\delta\), choose a binary
sequence \(a_{t,b_\beta}\in2^\omega\) representing \(s_{t,\beta}\).
The tuples and coordinates form a set, so these simultaneous choices
are available in ZFC.

The order \(B\) has the properties required by Lemma~\ref{aec:arrays}.
Apply that lemma to this family of sequences, using disjoint tagged
copies of \(I\), \(B\), and the other carriers as in its construction.
It gives a basic structure \(M_{\mathrm u}\) and a row-choice system
\(\mathbf z\) such that
\[
 c_t^{\mathbf z}(b_\beta)
   =[a_{t,b_\beta}]_{=^*}
   =s_{t,\beta}.
\]
If \(X\) is the underlying set of \(t\), Lemma~\ref{end:diagrams}
therefore makes its full subsequence test pass whenever
\(\beta\notin F_X\). In the notation of~\eqref{aec:bad-set},
\[
 E_t(\mathbf z)\subseteq\{b_\beta:\beta\in F_X\}.
\]
The right-hand side is countable. This verifies the object condition
for every tuple, so \(M_{\mathrm u}\in K\). The chosen representatives
verify that condition; they are not named in the structure.

\begin{theorem}\label{end:two-models}
The class \(K\) has two nonisomorphic models of
cardinality \(\beth_\eta\).  One has an order sort
of cardinality \(\aleph_1\), and the other has a
countable order sort.
\end{theorem}

\begin{proof}
The construction above and Lemma~\ref{end:hierarchy-size} give
\(M_{\mathrm u}\in K\) with
\[
  |I^{M_{\mathrm u}}|=|V_\eta|=\beth_\eta,
  \qquad
  |B^{M_{\mathrm u}}|=\aleph_1.
\]
Lemma~\ref{aec:size} gives total cardinality
\(\kappa=\beth_\eta\).

For \(M_{\mathrm c}\), take an index set of cardinality \(\kappa\),
base order \(B=\Q\), and constant zero sequences \(a_{t,b}\).
Lemma~\ref{aec:arrays} supplies a basic structure and makes the object
condition automatic because \(B\) is countable. Hence
\(M_{\mathrm c}\in K\), and Lemma~\ref{aec:size} gives
\(\|M_{\mathrm c}\|=\kappa\).

An isomorphism preserves the unary predicate
for the order sort.  It therefore cannot identify
\(M_{\mathrm u}\), whose order sort is uncountable,
with \(M_{\mathrm c}\), whose order sort is
countable.
\end{proof}
